\section{Technical Integration}

\subsection{Hybrid Shiny Architecture}
The dashboard uses a hybrid architecture that combines Python Shiny, React, and Plotly. Python Shiny acts as the application and deployment layer: at startup, \texttt{app.py} loads the cleaned EV and socioeconomic CSV files with Pandas, converts missing values and NumPy types into JSON-safe records, and injects the resulting data into the page before any visualization bundle is executed. React manages the dashboard interface and shared application state, while Plotly renders the charts and provides hover, zoom, pan, and click interactions. This separation preserves the layout flexibility of a client-side dashboard while retaining a reproducible Python entry point that can run locally or be deployed directly to shinyapps.io.

The JavaScript bundles are loaded in a deliberate dependency order. First, Shiny injects the cleaned CSV records; second, baseline IEA fixtures define the expected chart structures; third, a CSV bridge maps the cleaned records onto the visualization variables; finally, the chart and interface bundles are loaded. The bridge only replaces a series when a valid mapping is available, which prevents a partially missing dataset from breaking the complete dashboard. Static React, Plotly, and geographic TopoJSON assets are vendored within the repository, allowing the application to remain functional without external content-delivery networks. Asset modification timestamps are appended to bundle URLs to invalidate stale browser caches after deployment.

\subsection{Linked Interaction Design}
The interface is organized into six narrative tabs: Overview, Adoption Trends, Infrastructure and Demand, Market Composition, Socioeconomic and Policy, and ML/Causality. A collapsible filter sidebar exposes year range, development group, country, and EV powertrain controls. React stores these selections in a shared context and broadcasts each update through a custom \texttt{dashboard-filter-change} event. Chart builders receive the current filter state and apply it where a compatible country, development-group, powertrain, or temporal series exists, enabling coordinated filtering without a full page refresh or repeated server requests. Active-filter chips make the current selection visible and allow users to remove individual filters or reset the dashboard.

The choropleth provides a stronger linked-view interaction. Clicking a country dispatches a \texttt{dashboard-set-country} event, updates the global country filter, zooms the map to the selected country, and highlights or filters compatible charts across other tabs. The animated country-ranking chart uses \texttt{Plotly.react} to update yearly values efficiently, while the remaining charts are regenerated when their relevant filter dependencies change. Tooltips expose exact values and contextual variables without permanently overcrowding the visual encodings. Each chart card also supports PNG export, copying its underlying plotted data, and fullscreen inspection, allowing users to move from overview-level comparison to detailed analysis.

\subsection{Performance, Reproducibility, and Trade-offs}
Because the cleaned datasets are relatively small and the interaction logic is primarily visual, the server intentionally performs data loading once rather than recomputing every chart reactively. Filtering and rendering then occur in the browser, producing fast responses and reducing shinyapps.io server traffic. Plotly resize handlers ensure that charts remain correctly fitted when the sidebar, tabs, or viewport dimensions change. The deployment is reproducible through the pinned Python dependencies in \texttt{requirements.txt}, the documented local run command, and a single Shiny entry point.

This design goes beyond a basic collection of static charts by integrating spatial selection, linked filtering, animation, analytical model outputs, and export tools into one workflow. Its main trade-off is that the complete cleaned dataset is serialized to the browser at startup. This approach is efficient at the current scale, but a substantially larger dataset would benefit from Shiny reactive endpoints or modular server-side queries that return only the requested subset. The current architecture therefore prioritizes immediate interaction and visual consistency while preserving a clear path toward server-side scaling.
